% =====================================================
% 题型：离散随机变量期望与对称性选择题 (q_stats_expectation_symmetry.tex)
% 知识板块：概率期望
% 主思维：特殊值验证方法
% 副思维：对称性、化归转化
% 错因标签：暴力枚举、忽视去点影响
% 讲义用途：开场检测
% =====================================================
% 难度：★★ (中档)
% =====================================================
\newcommand{\qStatsExpectationSymmetryStem}{设 \(U = \left\{(x_1, x_2, x_3) \mid x_i \in \{-2, -1, 1, 2\}, i=1, 2, 3\right\}\) 为空间中 \(64\) 个点构成的集合。点 \(P(1, 1, 1)\)，记样本空间 \(\Omega = U \setminus \{P\}\)。从 \(\Omega\) 中随机取一个点。定义随机变量 \(X\) 如下：对 \(\Omega\) 中的每个点 \(A(x_1, x_2, x_3)\)，令 \(X(A) = x_1 + x_2 + x_3\)。则 \(X\) 的数学期望为 \hfill （\quad）}

\newcommand{\qStatsExpectationSymmetryOptions}{%
  \mcoptions[2]{\(-\dfrac{1}{21}\)}{\(-\dfrac{1}{63}\)}{\(0\)}{\(\dfrac{1}{7}\)}%
}

\newcommand{\qStatsExpectationSymmetryAnswer}{A}

\newcommand{\qStatsExpectationSymmetrySolution}{%
  【解题思路】\\
  利用坐标对称性求数学期望。首先，分析全集 \(U\) 中所有坐标的取值特征，利用对称性求得在全集 \(U\) 上变量 \(X = x_1 + x_2 + x_3\) 的总和为 \(0\)；然后，由于样本空间 \(\Omega\) 只是去掉了点 \(P(1, 1, 1)\)，故 \(\Omega\) 中所有点的变量 \(X\) 之和等于 \(U\) 中总和减去 \(X(P)\)；最后，用这个总和除以样本空间中的点数 \(63\)，即可求得期望值。\\
  【解析计算】\\
  集合 \(U = \left\{(x_1, x_2, x_3) \mid x_i \in \{-2, -1, 1, 2\}, i=1, 2, 3\right\}\) 中共有 \(4^3 = 64\) 个点。\\
  对每个点定义随机变量 \(X = x_1 + x_2 + x_3\)。\\
  在整个集合 \(U\) 中，每个坐标分量 \(x_i\) 的取值都是等可能地为 \(-2, -1, 1, 2\)。\\
  因为这四个数关于原点对称，其平均值为 \(\dfrac{-2 - 1 + 1 + 2}{4} = 0\)。\\
  所以在整个全集 \(U\) 中，变量 \(X\) 的总和为 \(0\)。\\
  样本空间 \(\Omega = U \setminus \{P\}\)，其中 \(P(1, 1, 1)\)。\\
  对于点 \(P\)，其变量值为 \(X(P) = 1 + 1 + 1 = 3\)。\\
  从全集总和中减去该点后，样本空间 \(\Omega\) 中所有点对应的变量 \(X\) 的总和为：\\
  \(\sum_{A \in \Omega} X(A) = 0 - X(P) = 0 - 3 = -3\)。\\
  样本空间 \(\Omega\) 中点的个数为 \(|\Omega| = 64 - 1 = 63\)。\\
  由于从 \(\Omega\) 中随机抽取一个点是等可能的，故 \(X\) 的数学期望为：\\
  \(E(X) = \dfrac{\sum_{A \in \Omega} X(A)}{|\Omega|} = \dfrac{-3}{63} = -\dfrac{1}{21}\)。\\
  故选 A。%
}

\newcommand{\qStatsExpectationSymmetryDiagnosis}{%
  部分学生拿到此题后，试图通过逐一罗列 63 个点的坐标并对其求和，或者去计算每个取值对应的概率分布列。这种硬算的方法极其耗时，且容易出错。关键在于识别并应用全集上的对称性。%
}

\newcommand{\qStatsExpectationSymmetryTeachingNotes}{%
  本题为一道巧妙的期望题。讲评时应向学生展示两种方法的对比：硬算分布列（耗时易错） vs 整体对称性思想（直接得出）。通过此题培养学生的“整体法”和“对称性”思维。%
}
